\section{Reproducibility, Scope, and Operational Rule}
\label{sec:discussion}

\subsection{RQ9: Can the Decision Chain Be Replayed?}
The anonymous artifact is self-contained and uses only declared dependencies. A clean environment imports no main-repository code, rebuilds the paper and supplement, and regenerates five main-paper figures and five data-linked tables. A fresh Faiss environment reproduces registered native action rows with zero top-$k$, recall, NDC, or action-order mismatch. The two prospective stages add 288,000 response cells and 112,000 raw timing measurements with zero top-$k$ mismatch. Role manifests verify zero overlap among selection, certification, and evaluation.

All records are committed with checksums and deterministic regeneration scripts; numeric claims are bound to repository records rather than generated prose. The anonymous artifact URL is supplied through the submission system.

\subsection{What the Evidence Establishes}
Budget response changes across registered graph rebuilds, data refresh, implementations, and scale. ICBA detects this change and converts a frozen candidate into a target-qualified completed decision. TCP recovers endpoint-relative work for recurring profiled queries and adds resolved value beyond target-global calibration on SIFT. Target-global calibration is the stronger default on Arxiv. Source one-rung can work, but the paired audit shows that a fixed action can match it or provide greater efficiency while both policies satisfy the risk contract.

The resulting maintenance rule has three steps. First, audit inherited policies against the target event and endpoint. Second, certify the least complex viable candidate and its fallback on disjoint target queries. Third, escalate from fixed to target-global to query-conditional policies only when held-out gain and lifecycle accounting justify the additional information.

\subsection{Boundaries}
The main conclusions are conditional on registered query populations, grids, SLAs, and build panels. The prospective panel has eight builds per dataset and one CPU configuration; the paired audit reuses that build panel for attribution. TCP assumes recurring profiled queries. The complete cold-history ledger is expressed in NDC, while actual timing covers serving search rather than truth acquisition, control, energy, or storage. DARTH, Ada-ef, Deep1M, and Vamana retain different native estimands and support scope rather than a common ranking. These boundaries define the next replication targets without changing the fixed-target decisions reported here.

\section{Conclusion}
Search-budget portability is an index-conditioned deployment decision. ICBA makes that decision auditable by separating candidate construction, target certification, fallback, held-out evaluation, and cost. The evidence rejects automatic reuse, confirms qualified recovery on new registered builds and query roles, and shows why strong simple baselines are essential: fixed-256 can match or outperform source one-rung under the common contract, target calibration provides the largest measured gains on the paired panel, and TCP is valuable when recurring query histories add measurable information. The practical outcome is a complexity-aware policy ladder backed by target evidence rather than a universal recovery rule.
